% !TeX root = ../../packet-and-principal-models.tex
\section{Scope and questions for further research}\label{sec:conclusion}
Under the execution premise, the fixed marked packet meets the normalized
curve at precisely three quotient points. The common lattice construction
and the unit twist give the stated principal models and six Weil points.
These conclusions keep the integral lattice, polarization and field marking
visible; they do not identify Mostaed's literal principal locus.

There are two natural continuations. First, determine whether the same
finite-domain method handles other norm generators and levels. Such an
extension needs a new arithmetic domain and new geometric certificates.
Second, identify the constructed principal family with a specified moduli
normalization, or prove that the integral data obstruct that identification.
The latter question must fix the lattice, level and polarization module.

The finite quotient family of Paper II starts from a primitive Fermat
factor, not from further points of this packet. Its Jacobian and descent
questions therefore do not constitute unfinished cases of the three-point
intersection theorem. The supporting Hodge note likewise concerns abstract
fixed-field types and does not add points to this curve.
